

\section{La Narration Interactive (TPRS) : Transformer l'input en intake}
space{0.3cm}
oindent

L'enseignement des langues anciennes, et singulièrement du grec ancien, se trouve aujourd'hui à la
croisée des chemins. Longtemps dominé par le paradigme de la Grammaire-Traduction, qui privilégie
l'analyse métalinguistique explicite au détriment de la compétence communicative, le champ
disciplinaire s'ouvre progressivement aux apports des Sciences de l'Acquisition des Langues Secondes
(SLA). Au cœur de cette mutation pédagogique se trouve la problématique centrale de la section
III.2.\cite{III-2-3-ref3} de cette thèse : le passage critique de l'\textbf{input} (les données
linguistiques brutes auxquelles l'apprenant est exposé) à l'\textbf{intake} (les données
effectivement traitées et intégrées par le système cognitif de l'apprenant).

Ce rapport de recherche se propose d'explorer en profondeur la méthodologie \textit{Teaching
Proficiency through Reading and Storytelling} (TPRS), ou Narration Interactive, comme vecteur
privilégié de cette transformation cognitive dans le cadre spécifique de l'enseignement du grec
ancien. Contrairement aux approches traditionnelles qui postulent que la connaissance des règles
(savoir déclaratif) précède la compétence (savoir procédural), le TPRS s'appuie sur le postulat
inverse : c'est l'exposition massive à un input compréhensible, engageant et répétitif qui génère la
compétence grammaticale implicite.\cite{III-2-3-ref1}

Cependant, l'application de cette méthode à une langue flexionnelle complexe comme le grec ancien,
caractérisée par une morphologie riche et une distance culturelle temporelle, soulève des défis
spécifiques. Comment saturer la mémoire de travail de l'apprenant avec des structures morphologiques
complexes sans provoquer de surcharge cognitive? Comment les mécanismes de répétition inhérents au
TPRS, tels que le \og Circling\fg{}, favorisent-ils la consolidation mnésique nécessaire à la
lecture fluide des textes classiques?

Pour répondre à ces questions, ce rapport mobilise une analyse interdisciplinaire croisant la
didactique des langues classiques (travaux de Bob Patrick, Lance Piantaggini, Christophe Rico), la
psycholinguistique (modèles de Bill VanPatten) et les neurosciences cognitives. Il s'articule autour
d'une démonstration rigoureuse de l'efficacité du TPRS à transformer l'input en intake, non par la
magie du récit, mais par une ingénierie précise de l'attention et de la mémoire.



\subsection{1\. cadre théorique : de la perception à l'intégration linguistique}

L'analyse de l'efficacité du TPRS nécessite d'abord de déconstruire les mécanismes cognitifs sous-
jacents à l'acquisition du langage. Il ne s'agit pas simplement d'exposer l'élève au grec, mais de
comprendre comment cet input devient une compétence.



\subsection{La dichotomie input/intake dans le modèle de vanpatten}

La distinction fondamentale opérée par les chercheurs en acquisition, notamment Bill VanPatten,
réside entre ce qui est disponible pour l'apprenant (\textit{input}) et ce qui est réellement traité
(\textit{intake}). Dans un cours de grec traditionnel, l'enseignant peut fournir un input riche en
expliquant une règle de syntaxe ou en lisant un texte de Xénophon. Cependant, si l'attention de
l'élève est focalisée sur le décodage lexical ou l'analyse métalinguistique, la forme grammaticale
(par exemple, la désinence de l'optatif) peut ne jamais être traitée comme une donnée linguistique
significative. Elle reste du \og bruit\fg{}.\cite{III-2-3-ref3}

VanPatten a formalisé ce processus dans sa théorie de l'\textbf{Input Processing (IP)}. Selon ce
modèle, les apprenants traitent l'input pour le sens avant de le traiter pour la forme.

\begin{itemize}\item \textbf{Priorité au Sens :} Si un apprenant de grec entend \textit{« Le philosophe (nominatif) frappe l'esclave (accusatif) »}, il s'appuiera d'abord sur l'ordre des mots ou la sémantique lexicale pour comprendre qui frappe qui, ignorant souvent les désinences casuelles si elles ne sont pas cruciales pour le sens immédiat.\cite{III-2-3-ref4}\item \textbf{Compétition Cognitive :} La mémoire de travail humaine ayant une capacité limitée, toute ressource allouée au décodage du vocabulaire inconnu ou à la traduction mentale est une ressource soustraite au traitement de la grammaire. C'est le « bandwidth problem » (problème de bande passante) évoqué par VanPatten.\cite{III-2-3-ref5}\end{itemize}Le TPRS intervient précisément ici comme une solution ergonomique pour le cerveau. En \og
abritant\fg{} le vocabulaire (c'est-à-dire en le limitant drastiquement et en assurant sa
compréhension parfaite avant la narration), le TPRS libère la mémoire de travail. L'élève, n'ayant
plus à lutter pour le sens lexical, peut enfin allouer son attention aux marques morphologiques
(l'intake grammatical).\cite{III-2-3-ref6}



\subsection{La saturation de la mémoire de travail et le rôle de l'attention}

Les neurosciences cognitives apportent un éclairage complémentaire sur la notion de saturation. Pour
qu'une structure linguistique passe de la mémoire de travail (traitement immédiat) à la mémoire à
long terme (stockage procédural), elle doit faire l'objet d'une répétition suffisante et d'une
attention focalisée.\cite{III-2-3-ref8}

Cependant, la saturation peut être à double tranchant :

1. \textbf{Saturation Négative (Surcharge) :} Si l'input est trop complexe, trop rapide ou trop
dense en informations nouvelles, la mémoire de travail sature (\og buffer overflow\fg{}).
L'apprenant décroche, le filtre affectif (anxiété) augmente, et l'apprentissage s'arrête. C'est
souvent le cas lors de la lecture traductionnelle de textes authentiques trop ardus pour le niveau
de l'élève.\cite{III-2-3-ref10}

2. \textbf{Saturation Positive (Consolidation) :} Pour acquérir une structure (ex: l'aoriste
sigmatique), le réseau neuronal correspondant doit être activé de manière répétée et intense sur une
courte période. Le TPRS, par ses techniques de questions-réponses cycliques, permet cette saturation
positive. Il bombarde le système cognitif avec la même structure cible sous des formes variées,
maximisant les chances que l'input devienne intake.\cite{III-2-3-ref12}



\subsection{L'acquisition implicite vs apprentissage explicite}

Stephen Krashen a longtemps soutenu que l'apprentissage conscient (règles de grammaire) ne peut
devenir de l'acquisition inconsciente (compétence fluide).\cite{III-2-3-ref3} Bien que cette
position soit parfois nuancée par d'autres chercheurs, le consensus en SLA reste que la compétence
communicative (la capacité de lire le grec couramment) repose sur des connaissances implicites,
procédurales.

Le TPRS vise directement la construction de cette compétence implicite. Contrairement à la méthode
Grammaire-Traduction qui sature la mémoire déclarative (règles, tableaux), le TPRS sature la mémoire
procédurale par l'action et la narration. Comme le souligne VanPatten, \og ce que vous pensez voir
dans la langue \[les règles\] n'est pas ce qu'est la langue dans l'esprit\fg{}.\cite{III-2-3-ref15}
La langue est un réseau complexe de connexions probabilistes que seul un input massif et
compréhensible peut construire.



\subsection{2\. anatomie méthodologique du tprs appliquée au grec}

Le TPRS n'est pas une méthodologie improvisée mais un protocole rigoureux en trois étapes, conçu
pour optimiser le traitement de l'input. Analysons comment chaque étape facilite spécifiquement
l'intake du grec ancien.



\subsection{Étape i : établir le sens (establish meaning)}

Cette étape prépare le terrain cognitif. Avant de pouvoir traiter grammaticalement une phrase
grecque, l'élève doit en posséder les briques sémantiques.

\begin{itemize}\item \textbf{Techniques d'Ancrage :} L'enseignant présente le nouveau vocabulaire (généralement limité à 3 structures, ex: \textit{lambanei} \- il prend, \textit{trechei} \- il court, \textit{lithos} \- la pierre).\item \textbf{Traduction et Gestuelle :} Contrairement à certaines méthodes directes puristes, le TPRS autorise la traduction rapide (en L1) pour établir le sens instantanément et éviter l'ambiguïté, facteur de stress cognitif. De plus, l'association d'un geste (issu du TPR classique de James Asher) permet un double encodage (sémantique et moteur), renforçant la trace mnésique.\cite{III-2-3-ref2}\item \textbf{Application au Grec :} Pour des concepts abstraits fréquents en philosophie ou tragédie grecque (ex: \textit{hubris}, \textit{dike}), l'enseignant peut utiliser des situations concrètes ou des analogies modernes pour ancrer le sens avant même que le mot ne soit inséré dans une phrase complexe.\cite{III-2-3-ref2}\end{itemize}

\subsection{Étape ii : la narration interactive (ask a story)}

C'est l'étape critique où l'input est transformé en intake grâce à la technique du
\textbf{Circling}. Loin d'être une simple répétition mécanique, le Circling est une exploration
systématique de la structure cible à travers des variations interrogatives.



\subsection{Le mécanisme du "circling" et la densité de l'input}

Prenons un exemple concret en grec ancien pour illustrer comment le Circling force le traitement
grammatical.

Phrase de base : Ho Socrates ten oikian zetei. (Socrate cherche la maison).

L'enseignant ne se contente pas de dire cette phrase. Il va \og cycler\fg{} autour d'elle pour
saturer l'audition de l'élève avec les désinences de l'accusatif (\textit{\-an}) et la forme verbale
(\textit{\-ei}), tout en vérifiant la compréhension du sens.\cite{III-2-3-ref17}

\begin{table}[h!]\small\centering\begin{tabular}{| p{2.3cm} | p{3.5cm} | p{3.5cm} |}\hline\textbf{Type de Question} & \textbf{Exemple en Grec (Translittéré)} & \textbf{Fonction Cognitive (VanPatten/Krashen)} \\ \hline\textbf{Affirmation} & \textit{Ho Socrates ten oikian zetei.} & Input initial ($i$). \\ \hline\textbf{Oui/Non (Positif)} & \textit{Ara ho Socrates ten oikian zetei?} (Oui) & Vérification basique du décodage lexical. \\ \hline\textbf{Oui/Non (Négatif)} & \textit{Ara ho Socrates ten hippon zetei?} (Non) & Discrimination sémantique (maison vs cheval). Force l'attention sur l'objet. \\ \hline\textbf{Alternative (Either/Or)} & \textit{Poteron ho Socrates ten oikian e ten hippon zetei?} & Choix forcé. L'élève doit produire le mot correct (Output minimal). \\ \hline\textbf{Question Sujet (Qui?)} & \textit{Tis ten oikian zetei?} & Focalisation sur le nominatif. \\ \hline\textbf{Question Objet (Quoi?)} & \textit{Ti zetei ho Socrates?} & Focalisation sur l'accusatif. \\ \hline\end{tabular}\end{table}Ce cycle permet à l'élève d'entendre la structure \textit{zetei} et l'accusatif \textit{ten oikian}
des dizaines de fois en quelques minutes, mais à chaque fois dans un contexte pragmatique différent
qui maintient l'éveil attentionnel. C'est cette \og inondation\fg{} (\textit{input flood}) contrôlée
qui permet au cerveau de commencer à segmenter le flux sonore et d'identifier les morphèmes
grammaticaux (Intake).\cite{III-2-3-ref4}



\subsection{La co-construction et l'engagement (pqa)}

Le TPRS se distingue par son aspect collaboratif. L'histoire n'est pas pré-écrite ; elle est \og
demandée\fg{}. L'enseignant pose des questions dont il n'a pas la réponse (ex: \og Socrate cherche
une maison... mais où? Athènes ou Sparte?\fg{}). Les élèves suggèrent des réponses.

\begin{itemize}\item \textbf{Personnalisation (PQA \- Personalized Questions and Answers) :} En intégrant les élèves ou leurs intérêts dans l'histoire, l'enseignant active le système limbique (émotion). Les neurosciences montrent que l'amygdale marque les souvenirs émotionnels comme prioritaires pour la consolidation à long terme.\cite{III-2-3-ref17}\item \textbf{Bizarre et Convaincant :} Les histoires TPRS tendent souvent vers l'absurde ou l'humour. Krashen parle de \og Compelling Input\fg{} (Input Irrésistible). Si l'histoire est suffisamment captivante, l'élève \og oublie\fg{} qu'il écoute une langue étrangère, abaissant le filtre affectif et maximisant l'intake inconscient.\cite{III-2-3-ref3}\end{itemize}

\subsection{Étape iii : lecture et discussion (read and discuss)}

Après l'oral, le passage à l'écrit est indispensable, surtout pour une langue de culture comme le
grec ancien.

\begin{itemize}\item \textbf{Transition Phonème-Graphème :} L'élève lit une version de l'histoire qu'il connaît déjà oralement. Cette familiarité réduit la charge cognitive de la lecture : l'élève ne déchiffre pas, il reconnaît. Cela permet de développer la fluidité de lecture (\textit{reading fluency}) bien plus rapidement que par la traduction mot-à-mot.\cite{III-2-3-ref2}\item \textbf{Grammaire Pop-up :} Lors de la lecture, l'enseignant peut faire de brèves incursions métalinguistiques (ex: \og Notez le \textit{nu} final ici, c'est l'objet\fg{}). Ces explications durent quelques secondes (\og pop-up\fg{}) et ne visent pas à enseigner la règle, mais à confirmer une intuition que l'élève a déjà développée grâce à l'input oral.\cite{III-2-3-ref5}\end{itemize}

\subsection{3\. adaptations spécifiques au grec ancien et innovations méthodologiques}

L'application du TPRS au grec ancien ne va pas sans défis. Les structures de la langue et la nature
des textes cibles imposent des adaptations que des institutions comme l'Institut Polis et des
enseignants-chercheurs ont développées.



\subsection{Le défi morphologique : tprs vs "living sequential expression" (lse)}

Le grec ancien présente une complexité morphologique (flexions) que l'anglais ou l'espagnol n'ont
pas. L'Institut Polis à Jérusalem, sous la direction de Christophe Rico, a développé une technique
complémentaire au TPRS : la \textbf{Living Sequential Expression (LSE)}.

Il est crucial de distinguer ces deux approches dans le cadre de la thèse, car elles travaillent
l'intake différemment.\cite{III-2-3-ref21}

\begin{table}[h!]\small\centering\begin{tabular}{| p{2.3cm} | p{3.5cm} | p{3.5cm} |}\hline\textbf{Caractéristique} & \textbf{TPRS (Teaching Proficiency through Reading and Storytelling)} & \textbf{LSE (Living Sequential Expression)} \\ \hline\textbf{Origine} & Blaine Ray (1990s), basé sur Krashen/Asher. & Christophe Rico (Polis), basé sur François Gouin (1880s). \\ \hline\textbf{Structure Narrative} & Intrigue avec conflit, personnages, dialogue. & Séquence logique d'actions chronologiques (Séries). \\ \hline\textbf{Objectif Cognitif} & Fluidité conversationnelle, vocabulaire varié, temps du récit. & Automatisation de séquences logiques, grammaire visuelle. \\ \hline\textbf{Exemple Grec} & \og Le philosophe veut une maison mais il n'a pas d'argent...\fg{} & \og Je me lève, je marche vers la porte, je saisis la poignée, j'ouvre...\fg{} \\ \hline\textbf{Traitement de l'Intake} & Variable, dépendant de l'interaction imprévisible. & Systématique, saturant une structure précise (ex: 1ère pers. sg.). \\ \hline\end{tabular}\end{table}L'innovation majeure pour le grec ancien est l'hybridation : utiliser le LSE pour installer
solidement les structures de base (saturer la mémoire procédurale avec les verbes de mouvement et
les cas concrets) et le TPRS pour développer la complexité syntaxique et la fluidité narrative. Le
LSE prépare l'intake grammatical en réduisant la charge cognitive (la logique de l'action guide la
phrase), ce qui permet ensuite au TPRS d'être plus efficace.\cite{III-2-3-ref21}



\subsection{La question de la "rigueur" et du contenu}

Une critique récurrente faite au TPRS dans les cercles classiques est le manque supposé de \og
rigueur\fg{} ou la pauvreté du contenu (histoires d'éléphants vs textes de Platon).

Cependant, les praticiens comme Bob Patrick et Lance Piantaggini démontrent que le TPRS permet en
réalité d'accéder plus vite aux textes authentiques.

\begin{itemize}\item \textbf{Novellas et \og Embedded Readings\fg{} :} Pour combler le fossé entre l'histoire TPRS simple et le texte antique, on utilise des lectures graduées (\textit{novellas}) ou des lectures \og enchâssées\fg{} (\textit{embedded readings}). On part d'une version très simplifiée (niveau TPRS) d'un mythe ou d'un épisode historique, et on ajoute progressivement de la complexité et du vocabulaire original, strate par strate.\cite{III-2-3-ref20}\item \textbf{Rigueur Cognitive vs Rigueur Philologique :} La véritable rigueur, argumentent ces chercheurs, n'est pas d'imposer une complexité indigeste dès le début (ce qui mène à l'échec ou à l'abandon), mais de respecter les contraintes neurocognitives de l'apprentissage pour mener l'élève à une véritable compétence de lecture.\cite{III-2-3-ref26}\end{itemize}

\subsection{Ressources et communauté de pratique}

Contrairement aux langues vivantes, le corpus de matériel TPRS grec est en construction.

\begin{itemize}\item \textbf{Lance Piantaggini (Magister P)} a théorisé l'usage des \textit{novellas} (petits romans en latin/grec facile) comme outils de saturation de l'input. Ses travaux montrent que la lecture extensive de ces textes simples consolide l'intake bien mieux que la lecture intensive de textes courts et difficiles.\cite{III-2-3-ref20}\item \textbf{Justin Slocum Bailey} propose des techniques pour générer de l'input compréhensible à partir de n'importe quel texte, transformant la lecture laborieuse en session de questions-réponses vivante (\og Est-ce que César a traversé le Rubicon ou le Tibre?\fg{}).\cite{III-2-3-ref24}\end{itemize}

\subsection{4\. analyse des données et preuves d'efficacité}

La validation de la section III.2.\cite{III-2-3-ref3} de votre thèse doit s'appuyer sur des données
probantes. Bien que les études longitudinales spécifiques au grec ancien soient encore rares,
l'extrapolation des données sur le TPRS en général et en latin permet de tirer des conclusions
solides.



\subsection{La synthèse de karen lichtman ("the wonder table")}

Karen Lichtman, chercheuse en SLA, a compilé et analysé une vaste quantité d'études comparatives sur
le TPRS. Son analyse, souvent désignée comme \og The Wonder Table\fg{}, offre des arguments
statistiques puissants.\cite{III-2-3-ref31}

\begin{itemize}\item \textbf{Dominance du TPRS :} Sur 29 études comparatives rigoureuses, 16 démontrent une supériorité statistiquement significative du TPRS sur les méthodes traditionnelles, tant en vocabulaire qu'en grammaire (intake). Seulement 4 études favorisent les méthodes non-TPRS, souvent dans des contextes de tests purement explicites (récitation de règles).\item \textbf{Facteur de Rétention :} Un point crucial pour les langues anciennes (où le taux d'abandon est élevé) est que le TPRS montre systématiquement de meilleurs résultats sur les facteurs affectifs : motivation, plaisir, et continuation des études.\cite{III-2-3-ref32}\end{itemize}

\subsection{Données neurocognitives sur la répétition et l'intake}

Les mécanismes du TPRS sont validés par la recherche fondamentale sur la mémoire.

\begin{itemize}\item \textbf{Consolidation Procédurale :} La mémoire procédurale (savoir-faire) nécessite une répétition massive pour se stabiliser. Les études sur la répétition lexicale montrent qu'il faut rencontrer un mot entre 10 et 50 fois en contexte pour qu'il soit acquis (intake). Le TPRS, par le Circling, permet d'atteindre ces fréquences en une seule séance, ce que la traduction linéaire ne permet jamais (un mot rare chez Homère peut n'apparaître qu'une fois).\cite{III-2-3-ref8}\item \textbf{Gestion de la Charge Cognitive :} Les études sur la saturation de la mémoire de travail confirment que lorsque la charge est trop élevée (trop d'informations nouvelles), l'apprentissage s'effondre. Le TPRS respecte l'architecture cognitive en limitant le débit d'information nouvelle pour maximiser le traitement de l'information présente.\cite{III-2-3-ref10}\end{itemize}

\subsection{Critiques et réponses}

Il est intellectuellement honnête d'aborder les limites. Certains critiques (souvent issus de la
tradition classique) arguent que le TPRS ne prépare pas assez à la complexité analytique des textes
littéraires.

\begin{itemize}\item \textbf{Réponse de VanPatten :} La compétence analytique (savoir décrire la langue) est distincte de la compétence linguistique (savoir utiliser la langue). Le but du TPRS est la compétence. L'analyse peut venir \textit{après}, une fois la langue acquise, mais elle ne peut la \textit{remplacer}.\cite{III-2-3-ref3}\item \textbf{Le Mythe de la Rigueur :} Comme le note Bob Patrick, la méthode traditionnelle, malgré sa \og rigueur\fg{} apparente, échoue pour la majorité des élèves (qui ne lisent jamais le grec couramment). Le TPRS propose une rigueur différente : celle de l'acquisition effective pour tous.\cite{III-2-3-ref27}\end{itemize}

\subsection{5\. synthèse bibliographique stratégique pour la thèse}

Cette section propose une organisation raisonnée des sources pour étayer le chapitre III.2.3, en
classant les références par fonction argumentative.



\subsection{Fondements épistémologiques (le "pourquoi")}

Ces références justifient le changement de paradigme vers l'Input Compréhensible.

\begin{itemize}\item \textbf{Krashen, S. D.} (1982, 2013) : Pour les concepts d'Input, de Filtre Affectif et la distinction Acquisition/Apprentissage.\cite{III-2-3-ref2}\item \textbf{VanPatten, B.} (1996, 2004, 2017) : Essentiel pour la théorie de l'Input Processing, la saturation de la mémoire de travail et le rôle de l'attention focalisée sur la forme (\textit{Processing Instruction}).\cite{III-2-3-ref4}\item \textbf{Asher, J.} (2009) : Pour le lien entre mouvement physique et rétention lexicale (TPR), précurseur du TPRS.\cite{III-2-3-ref2}\end{itemize}

\subsection{Méthodologie et pratique (le "comment")}

Ces sources décrivent les mécanismes techniques (Circling, Story-asking).

\begin{itemize}\item \textbf{Ray, B. \& Seely, C.} (2015) : \textit{Fluency Through TPR Storytelling}. Le manuel de référence pour la mise en œuvre technique.\cite{III-2-3-ref1}\item \textbf{Lichtman, K.} (2014, 2018) : Pour les méta-analyses sur l'efficacité du TPRS (\textit{The Wonder Table}) et les études comparatives.\cite{III-2-3-ref31}\item \textbf{Piantaggini, L.} (Magister P) : Pour l'adaptation des techniques de narration et la création de \textit{novellas} en langues classiques.\cite{III-2-3-ref20}\end{itemize}

\subsection{Application au grec et latin (le contexte spécifique)}

Ces sources ancrent la recherche dans la didactique des langues anciennes.

\begin{itemize}\item \textbf{Patrick, R.} (2019) : \textit{Comprehensible Input and Krashen's Theory} (Journal of Classics Teaching). Un article fondateur pour la communauté classique.\cite{III-2-3-ref17}\item \textbf{Rico, C. \& Institut Polis} (2015, 2021) : Pour la méthode Polis, le \textit{Living Sequential Expression} (LSE) et l'enseignement du grec Koinè comme langue vivante.\cite{III-2-3-ref21}\item \textbf{Slocum Bailey, J.} : Pour les techniques de questions-réponses appliquées aux textes antiques (\textit{Indwelling Language}).\cite{III-2-3-ref30}\end{itemize}

\subsection{Conclusion}

L'analyse approfondie de la section III.2.\cite{III-2-3-ref3} permet de conclure que la Narration
Interactive (TPRS) constitue un levier pédagogique puissant et théoriquement fondé pour
l'enseignement du grec ancien. En respectant les contraintes neurobiologiques de la mémoire de
travail et en exploitant les mécanismes de l'Input Processing décrits par VanPatten, le TPRS
parvient à transformer l'input éphémère en intake durable.

Loin d'être une dilution du contenu classique, l'approche par histoires répétitives et questions-
réponses offre une rigueur cognitive supérieure à la méthode traditionnelle. Elle permet de
construire, strate par strate, la compétence implicite nécessaire à la lecture fluide. L'intégration
des innovations spécifiques au grec, telles que le \textit{Living Sequential Expression} de
l'Institut Polis, enrichit encore ce modèle en offrant des outils adaptés à la morphologie complexe
de la langue.

Ainsi, dans le cadre d'une thèse sur l'enseignement du grec, le TPRS ne doit pas être présenté comme
une simple \og activité ludique\fg{}, mais comme une \textbf{technologie de l'acquisition} validée
par la recherche, capable de répondre à la crise de l'enseignement des langues anciennes en rendant
le grec accessible, vivant et, in fine, acquis.

Ce tableau synthétise les différences techniques entre les approches pour faciliter l'analyse dans
la thèse.

\begin{table}[h!]\small\centering\begin{tabular}{| p{2cm} | p{2.2cm} | p{2.2cm} | p{2.2cm} |}\hline\textbf{Variable Didactique} & \textbf{Grammaire-Traduction (Traditionnel)} & \textbf{TPRS (Ray / Patrick)} & \textbf{LSE (Polis / Rico)} \\ \hline\textbf{Unité d'Input} & La phrase écrite, isolée ou dans un texte littéraire dense. & L'histoire (récit) co-construite oralement. & La séquence d'actions logiques (Série) orale et physique. \\ \hline\textbf{Rôle de la Grammaire} & Objet d'étude explicite (Savoir). & Outil de communication implicite (Compétence). \og Pop-up grammar\fg{}. & Structure induite par la logique de l'action (\og Je... tu... il...\fg{}). \\ \hline\textbf{Gestion de la Mémoire} & Surcharge fréquente (Lexique \+ Grammaire simultanés). & Optimisation : Lexique abrité, Grammaire par répétition (Circling). & Optimisation : Séquence logique aide la mémorisation (Script cognitif). \\ \hline\textbf{Validation de l'Intake} & Traduction correcte (Output différé). & Réponse immédiate aux questions (PQA), fluidité. & Exécution correcte de l'action ou description de la série. \\ \hline\textbf{Adaptation au Grec} & Totale pour l'analyse textuelle, nulle pour l'oral. & Forte pour l'acquisition des cas par l'oreille. & Très forte pour l'automatisation des verbes et prépositions. \\ \hline\end{tabular}\end{table}